% paper/sections/problem_formulation.tex
% §3.1 Babaei (2013) unconstrained formulation
% §3.2 Our budget-constrained extension

\section{Problem Formulation}
\label{sec:problem}

We formalize two variants of the revenue maximization problem:
the unconstrained variant of \citet{babaei2013revenue} and our
budget-constrained extension.

\subsection{Unconstrained Revenue Maximization}
\label{sec:babaei}

Revenue maximization exploits the fact that social influence raises a buyer's
willingness-to-pay.

\begin{definition}[Influence Network]
Let $G = (\calV, \calE)$ be a directed weighted graph with edge weights
$w_{ij} \in [0,2]$ drawn independently from $\mathrm{Uniform}(0,2)$,
representing the influence strength of buyer $j$ on buyer $i$.
\end{definition}

\begin{definition}[Valuation]
Given seed set $\St \subseteq \calV$, the \emph{normalized influence} on buyer
$i$ is
\[
  x_i(\St) = \frac{\sum_{j \in \St \cup \{i\}} w_{ij}}{\sum_{k \in \calV} w_{ik}},
\]
and the \emph{valuation} (willingness-to-pay) of buyer $i$ is
$\val_i(\St) = f\!\left(2 x_i(\St)\right)$,
where $f$ is the Rayleigh CDF: $f(y) = (y/b^2)\exp(-y^2/2b^2)$, $b=1$.
The monotone variant clips $f$ to be non-decreasing (used for Forest-Fire graphs);
the non-monotone variant uses the raw Rayleigh PDF (used for Rice-Facebook).
\end{definition}

\begin{definition}[Offer and Acceptance]
At step $t$ the seller presents buyer $v^* \notin \St$ with price
$p_t = \val_{v^*}(\St) \cdot (1 - d_t)$, where $d_t \in [0,1]$ is the discount.
Buyer $v^*$ accepts if $p_t \le \val_{v^*}(\St)$
(equivalently, always accepts since $p_t \le \val_{v^*}(\St)$ by construction).
In practice we model stochastic acceptance: $v^*$ accepts with probability
$p_t / \val_{v^*}(\St) = 1 - d_t$ when the seller does not know the exact
weight realization and uses 200 Monte Carlo samples to estimate
$\hat{\val}_{v^*}(\St)$.
\end{definition}

Revenue from step $t$ is $r_t = p_t \cdot \mathbf{1}[\text{accepted}]$.
The seller's objective is:

\begin{problem}[Unconstrained Revenue Maximization]\label{prob:uncon}
Choose a sequence of buyers $v_1, \ldots, v_n$ and discounts $d_1, \ldots, d_n$
to maximize total revenue $R = \sum_{t=1}^{n} r_t$,
over all $n = |\calV|$ sequential offers.
\end{problem}

\noindent
The key difficulty is the \emph{joint coupling}: the revenue-optimal discount for
$v^*$ at step $t$ depends on the current influence $x_{v^*}(\St)$, which depends on
which buyers accepted in steps $1, \ldots, t-1$.
\citet{babaei2013revenue} ignore this coupling; we learn it end-to-end.

\subsection{Budget-Constrained Revenue Maximization}
\label{sec:budget}

Budget constraints transform seed selection into a sequential knapsack with
network externalities.

\begin{definition}[Production Cost and Capital]
The seller starts with capital $B_0 = k \cdot c$ for a given
\emph{initial target count} $k$ and unit production cost $c > 0$.
Each offer costs $c$ regardless of acceptance; successful sales replenish the
budget by the accepted price $p_t$.
The budget after step $t$ is:
\[
  B_{t} = B_{t-1} - c + r_t.
\]
\end{definition}

\begin{definition}[Bankruptcy]
The seller goes bankrupt at the first step $t$ for which $B_t < c$.
Once bankrupt, no further offers can be made.
\end{definition}

\begin{problem}[Budget-Constrained Revenue Maximization]\label{prob:budget}
Maximize total revenue $R = \sum_t r_t$ subject to solvency: $B_t \ge 0$ for all $t$.
\end{problem}

\noindent
Problem~\ref{prob:budget} strictly generalizes Problem~\ref{prob:uncon}:
as $c \to 0$ the constraint becomes vacuous.
With finite $c$, the policy must balance two competing effects.
\begin{itemize}[leftmargin=*,topsep=2pt,itemsep=0pt]
  \item \emph{Deep discounts} increase acceptance probability and revenue per step,
    replenishing the budget for future offers.
  \item \emph{Shallow discounts} (high prices) yield larger revenue per accepted
    buyer but risk rejection, draining the budget without replenishment.
\end{itemize}
Skipping low-value buyers (those unlikely to accept at a profitable price) becomes
a legitimate strategy to preserve budget for high-value opportunities later —
a decision that is NP-hard to optimize globally~\citep{papadimitriou1994complexity}.

\paragraph{Regime characterization.}
Let $k$ be the initial target count and $c = 0.3$ (our experimental setting).
The budget $B_0 = kc$ spans three qualitatively different regimes:
\begin{enumerate}[leftmargin=*,topsep=2pt,itemsep=0pt]
  \item \emph{Capital-binding} ($k \le 5$, $B_0 \le 1.5$):
    a single rejection can cause bankruptcy; skip decisions dominate.
  \item \emph{Transitional} ($6 \le k \le 15$, $1.8 \le B_0 \le 4.5$):
    skipping is beneficial but not mandatory; rich solution landscape.
  \item \emph{Budget-slack} ($k \ge 20$, $B_0 \ge 6.0$):
    capital rarely binds; the problem reduces approximately to the unconstrained
    variant.
\end{enumerate}
Our algorithms are evaluated across all three regimes (§\ref{sec:exp_budget}).
